\section{Hierarchical plans and reusable skills}
\label{sec:C10:hierarchy}
Hierarchical planning replaces some primitive actions with temporally extended procedures. An option can be described by an initiation set, an internal policy, and a termination condition. For almost-surely finite duration $\tau\geq1$, bounded rewards and discount $0\leq\gamma<1$, a semi-Markov value expression is $Q(s,o)=\mathbb E[\sum_{i=0}^{\tau-1}\gamma^ir_i+\gamma^\tau\max_{o'}Q(s',o')]$. The discount exponent records elapsed primitive time or steps under the chosen convention. Treating every skill as one equal-cost action ignores duration, tool expenditure, and failure exposure. A skill library can therefore make a plan shorter on paper while leaving its actual operational horizon large.

Reusable procedures need applicability conditions. A code-editing skill can depend on repository layout, language version, permissions, or a particular test command. A browser skill can depend on an interface that later changes. The planner should validate these conditions before invocation and interpret termination based on observed state. A natural-language skill description can help the model select a procedure, but it does not establish that the installed implementation satisfies the description. Versioned executable resources and tests provide a stronger connection between symbolic planning and actual execution.

Abstraction can also conceal conflicts between subgoals. Two procedures may independently succeed while writing incompatible state, consuming the same budget, or invalidating each other's assumptions. A dependency graph should include shared resources and state versions, not merely an ordered list of textual tasks. Parallel execution is safe only when the relevant effects commute or are coordinated through an appropriate transaction or synchronization mechanism. Assigning different agent roles does not make their writes independent. This is a planning issue before it becomes a distributed-runtime issue.

A useful hierarchy experiment compares primitive control, fixed expert-authored skills, learned skill retrieval, and adaptive skill revision under the same external interfaces. Measure initiation errors, termination errors, task success, and complete cost. Hold out tool or repository variants to test transfer beyond memorized workarounds. A shorter generated plan is not by itself an improvement; it may simply hide more work inside a skill call. The scientific question is whether the abstraction improves decision quality and resource use while preserving the required constraints.

